% =====================================================================
% supplementary.tex --- Appendices for the EMNLP 2026 submission
% "Dose-Response Curves for Construction Learning in Small Language Models"
%
% Compiles standalone. To bind these as appendices of main.tex instead,
% move the \appendix block and \section{...} contents into main.tex and
% delete the preamble here. See main.tex for notes on obtaining acl.sty.
% =====================================================================

\IfFileExists{acl.sty}{%
  \documentclass[11pt]{article}
  \usepackage[review]{acl}
}{%
  \documentclass[11pt]{article}
  \usepackage[margin=1in]{geometry}
  \usepackage{natbib}
  \bibliographystyle{plainnat}
  \PackageWarningNoLine{supp}{acl.sty not found --- using plain article
    fallback. Install the EMNLP 2026 style before submitting.}
}

\usepackage[T1]{fontenc}
\usepackage[utf8]{inputenc}
\usepackage{graphicx}
\usepackage{booktabs}
\usepackage{amsmath}
\usepackage{xcolor}
\usepackage{url}

\newcommand{\resultpending}[1]{%
  \textcolor{red}{\textbf{[RESULT PENDING: #1]}}}
\newcommand{\TODO}[1]{\textcolor{blue}{\textbf{[TODO: #1]}}}

\title{Supplementary Material:\\
  Dose-Response Curves for Construction Learning in\\
  Small Language Models}
\author{Anonymous Submission}

\begin{document}
\maketitle
\appendix

% =====================================================================
\section{Filter Specifications}
\label{app:filters}
Each construction is detected by a dedicated filter over Stanza dependency
parses \citep{qi2020stanza}, with surface fallbacks where parses are
brittle. We re-run the relevant filter after building each dose corpus to
verify realized counts.

\paragraph{AANN (article + adjective + numeral + plural noun).}
We match the four-token core \emph{article ADJ NUM Noun.Plural} (as in
``a beautiful five days'') and confirm the article actually depends on the
plural noun head, which rules out coincidental sequences that span a clause
boundary. The semantics treat the plural as a measured quantity, which is
why an indefinite article can scope over a plural NP
\citep{mahowald2023aann}.

\paragraph{Comparative correlative (``the X-er $\ldots$, the Y-er $\ldots$'').}
We detect two ``the + comparative'' openings split by a comma, with the
first clause at the start of the sentence. The ``the'' here is a frozen
degree marker, not an ordinary article. A regex fallback over raw text
recovers cases where the parser mangles the unusual syntax (e.g.\ ``the
more the merrier''), boosting recall \citep{weissweiler2022comparative}.

\paragraph{Tough-movement (``this book is tough to read'').}
We match the reliable surface frame \emph{BE + tough-adjective + ``to'' +
VERB} using a fixed tough-class adjective lexicon. We deliberately do not
verify the syntactic gap, which is brittle to parse; the lexicalized
surface frame is precise enough \citep{hu2020syntaxgym}.

\paragraph{Resultative (V + NP + AP naming a result state).}
We require, in the dependency parse, a verb that governs both an \texttt{obj}
and an adjectival \texttt{xcomp} (the result phrase), as in ``hammered the
metal flat.'' Testing the dependency structure directly avoids reasoning
about result semantics from the surface string
\citep{goldberg2004resultatives}.

\paragraph{Dose corpus construction.}
Removed positive sentences are replaced by matched sentences from a
100M-word pool, in priority order: (1) same source domain, (2) length
within $\pm 20\%$, relaxing to $\pm 40\%$, then (3) nearest-length
sentence from the same domain. Corpus generation uses a fixed seed (42),
and every corpus ships with an audit JSON of kept/removed/swapped
sentences. \resultpending{realized attested-``all'' counts per construction;
fill from results/dose\_corpora\_sanity.json}

% =====================================================================
\section{Evaluation Item Schema and Counts}
\label{app:items}
Evaluation items are stored as JSONL minimal pairs, one per line, with the
fields below.

\begin{table}[h]
  \centering
  \small
  \begin{tabular}{ll}
    \toprule
    Field & Description \\
    \midrule
    \texttt{item\_id}          & unique identifier \\
    \texttt{construction}      & target construction \\
    \texttt{good\_sentence}    & grammatical member \\
    \texttt{bad\_sentence}     & minimally altered ungrammatical member \\
    \texttt{minimal\_pair\_type} & contrast type (e.g.\ word\_order) \\
    \texttt{source}            & provenance / template note \\
    \texttt{notes}             & rationale for the contrast \\
    \bottomrule
  \end{tabular}
  \caption{Evaluation item schema (JSONL).}
  \label{tab:item-schema}
\end{table}

A model is scored correct on a pair when
$\mathrm{SLOR}(\text{good}) > \mathrm{SLOR}(\text{bad})$. Per-construction
item counts:

\begin{table}[h]
  \centering
  \small
  \begin{tabular}{lc}
    \toprule
    Construction & \# minimal pairs \\
    \midrule
    AANN                    & 30 \\
    Comparative correlative & \resultpending{count once items finalized} \\
    Tough-movement          & \resultpending{count once items finalized} \\
    Resultative             & \resultpending{count once items finalized} \\
    \bottomrule
  \end{tabular}
  \caption{Evaluation item counts. AANN items are finalized (30 pairs);
    the remaining sets are in progress (\texttt{evaluation/eval\_items/}).}
  \label{tab:item-counts}
\end{table}

% =====================================================================
\section{Training Hyperparameters}
\label{app:hyperparams}
All values below are frozen across all runs
(\texttt{configs/base.yaml}). They are the BabyLM strict-track winning
configuration for LTG-BERT \citep{samuel2023ltgbert}; we do not tune them
per construction or dose.

\begin{table}[h]
  \centering
  \small
  \begin{tabular}{ll}
    \toprule
    Hyperparameter & Value \\
    \midrule
    \multicolumn{2}{l}{\emph{Architecture (LTG-BERT base)}} \\
    Hidden size              & 384 \\
    Layers                   & 12 \\
    Attention heads          & 6 \\
    Intermediate size        & 1024 \\
    Max position embeddings  & 128 \\
    Vocabulary size          & 16{,}384 \\
    \midrule
    \multicolumn{2}{l}{\emph{Optimization}} \\
    Optimizer                & AdamW \\
    $\beta_1, \beta_2$       & 0.9, 0.98 \\
    $\epsilon$               & $1\times10^{-6}$ \\
    Weight decay             & 0.1 \\
    Peak learning rate       & $1\times10^{-3}$ \\
    LR schedule              & cosine (decay to 0) \\
    Warmup ratio             & 0.1 \\
    Max grad norm            & 1.0 \\
    \midrule
    \multicolumn{2}{l}{\emph{Training}} \\
    Batch size               & 256 sequences \\
    Max sequence length      & 128 \\
    Epochs                   & 20 \\
    Precision                & bf16 \\
    MLM probability          & 0.15 (whole-word) \\
    Held-out slice           & 100k words \\
    \bottomrule
  \end{tabular}
  \caption{Frozen training configuration (\texttt{configs/base.yaml}).}
  \label{tab:hyperparams}
\end{table}

% =====================================================================
\section{Reproducibility Checklist}
\label{app:repro}
\begin{itemize}
  \item \textbf{Data.} BabyLM-10M corpus and 100M-word replacement pool,
        both public; parsed with Stanza \citep{qi2020stanza}.
  \item \textbf{Dose corpora.} Generated once with corpus seed 42 and
        frozen; per-corpus audit JSON records every kept/removed/swapped
        sentence; filters re-run to verify counts.
  \item \textbf{Models.} LTG-BERT trained from scratch, frozen recipe
        (Appendix~\ref{app:hyperparams}); training seeds 42, 43, 44; 60
        runs total. Implementation: Hugging Face Transformers
        \citep{wolf2020transformers}.
  \item \textbf{Evaluation.} SLOR with masked-LM pseudo-log-likelihood
        \citep{pauls2012slor}; unigram correction from each model's own
        dose corpus; every model scored on all four test sets; $n$-gram
        baseline as control. Results streamed row-by-row with resume.
  \item \textbf{Analysis.} Hill fits via bounded NLS with parametric
        bootstrap CIs (1{,}000 resamples; \citealp{efron1993bootstrap});
        AIC/BIC model comparison; $k$-means clustering with silhouette.
  \item \textbf{Decision.} Pre-registered rule applied by
        \texttt{drc.analysis.decision}; verdict written to
        \texttt{results/decision.txt}.
  \item \textbf{Artifacts released.} Code, frozen corpora, evaluation
        items, fitted curves, and figures.
  \item \textbf{Compute.} \resultpending{report total GPU hours and
        hardware once runs complete}
\end{itemize}

% =====================================================================
\section{Pre-Registered Hypotheses and Decision Rules}
\label{app:hypotheses}
We locked the following five hypotheses, their thresholds, and the
hypothesis-to-title mapping before inspecting any fitted numbers
(\texttt{src/drc/analysis/decision.py}). The rules are checked in order;
if more than one fires, or none fires cleanly, we report ``Mixed'' and
point back to the Outcome Matrix rather than forcing a headline.

\begin{itemize}
  \item \textbf{H1 --- Smooth scaling.} Hill wins by AIC for every
        construction \emph{and} mean Hill coefficient $n \in [0.5, 2.0]$.
        Learning is a gentle saturating climb.
        \emph{Title:} ``A Scaling Law for Construction Learning in
        Language Models.''
  \item \textbf{H2 --- Phase transition.} Hill wins everywhere but mean
        $n > 4$. Same curve, far steeper: grammatical grokking.
        \emph{Title:} ``Grammatical Grokking: Phase Transitions in
        Construction Acquisition.''
  \item \textbf{H3 --- Heterogeneous regimes.} $k{=}2$ clusters separate
        cleanly (silhouette $> 0.5$), or Hill wins for some constructions
        while the step model wins for others. Two routes to grammar.
        \emph{Title:} ``Two Routes to Grammar: Heterogeneous Dose-Response
        in Construction Learning.''
  \item \textbf{H4 --- Indirect evidence dominates.} The floor-to-ceiling
        gain is tiny: mean $(E_{\max}-E_0)/E_{\max} < 0.1$. Direct exposure
        barely moves the needle.
        \emph{Title:} ``Direct Exposure Doesn't Matter: Indirect Evidence
        in Construction Learning.''
  \item \textbf{H5 --- Stochastic acquisition.} Mean within-cell CV across
        seeds exceeds mean between-cell CV across doses. Seed noise swamps
        the dose signal.
        \emph{Title:} ``Stochastic Construction Acquisition: When Seeds
        Beat Doses.''
\end{itemize}

\paragraph{Verdict.}
\resultpending{which hypothesis fired and the supporting statistics; fill
from results/decision.txt}

% =====================================================================
\bibliography{references}

\end{document}
